% 07_geophysics.tex
% Ch.7: Seismic PREM on the FDTD engine

\chapter{Geophysics: Seismic Waves}
\label{ch:geophysics}

\section{Seismic Wave Propagation on the FDTD Engine}
\label{sec:seismic_fdtd_bridge}

The 3D FDTD Maxwell solver (\texttt{fdtd\_3d.py}) is a \textit{universal}
wave engine.  Injecting the PREM Earth model's impedance profile as
material maps ($\varepsilon_r$, $\mu_r$) along the radial axis
produces correct seismic wave propagation, reflection coefficients
at layer boundaries, and waveguide trapping in the low-velocity zone.

\section{Constitutive Mapping}

\begin{align}
    \varepsilon_r(r) &= K_{\text{ref}} / K(r)
    \quad\text{(compressibility $\to$ capacitance)} \\
    \mu_r(r) &= G_{\text{ref}} / G(r)
    \quad\text{(shear compliance $\to$ inductance)}
\end{align}
and the reflection coefficient at each boundary is:
\begin{resultbox}{Seismic Reflection Coefficient (Moho)}
\begin{equation}
    \Gamma_{\text{Moho}} = \frac{\rho_2 V_{p2} - \rho_1 V_{p1}}
    {\rho_2 V_{p2} + \rho_1 V_{p1}}
\end{equation}
\end{resultbox}
which is the \textit{same} \texttt{reflection\_coefficient(Z1, Z2)}
used for Pauli exclusion, plasma cutoff, superconductor boundaries,
and antenna port $S_{11}$.

\begin{figure}[H]
    \centering
    \includegraphics[width=0.95\textwidth]{seismic_prem_impedance.pdf}
    \caption{PREM Earth model on the AVE FDTD engine.
    \textbf{Left:} Compressional ($V_P$) and shear ($V_S$) velocity
    profiles through the Earth (surface to inner core).
    \textbf{Centre:} Corresponding seismic impedance $Z = \rho V$.
    \textbf{Right:} Reflection coefficients at each layer boundary,
    computed by the same \texttt{reflection\_coefficient()} function
    used for all other domains.}
    \label{fig:seismic_prem}
\end{figure}

\section{PREM Layer Data and Numerical Evaluation}

Table~\ref{tab:prem_layers} lists the major discontinuities of the Preliminary Reference Earth Model (PREM) and the corresponding seismic impedance computed by the AVE engine.

\begin{table}[H]
\centering
\caption{PREM layer impedances. Reflection coefficients are computed by the universal \texttt{reflection\_coefficient(Z1, Z2)} function---identical to the function used for particle boundaries, plasma cutoffs, and antenna ports.}
\label{tab:prem_layers}
\small
\begin{tabular}{@{}l r r r r@{}}
\toprule
\textbf{Boundary} & $\rho$ (kg/m$^3$) & $V_P$ (km/s) & $Z = \rho V_P$ (MPa$\cdot$s/m) & $|\Gamma|$ \\
\midrule
Surface (crust)       & 2,600  & 5.8   & 15.1 & --- \\
Moho (crust/mantle)   & 3,380  & 8.1   & 27.4 & 0.29 \\
670 km (upper/lower)  & 3,990  & 10.3  & 41.1 & 0.20 \\
CMB (mantle/core)     & 9,900  & 8.1   & 80.2 & 0.32 \\
ICB (outer/inner)     & 12,760 & 11.0  & 140.4 & 0.27 \\
\bottomrule
\end{tabular}
\end{table}

\noindent The Moho reflection coefficient evaluates numerically:
\begin{equation}
    \Gamma_{\text{Moho}} = \frac{Z_{\text{mantle}} - Z_{\text{crust}}}{Z_{\text{mantle}} + Z_{\text{crust}}} = \frac{27.4 - 15.1}{27.4 + 15.1} = \frac{12.3}{42.5} \approx 0.29
\end{equation}
This is the \textit{same} algebraic formula that produces Pauli exclusion ($\Gamma = -1$) at the saturated particle boundary and gravitational stealth ($\Gamma = 0$) in a gravity well. The only difference is the input impedances.

\section{Waveguide Trapping in the Low-Velocity Zone}

Between 80 and 220~km depth, the PREM model exhibits a pronounced low-velocity zone (LVZ) where partial melting reduces $V_S$ by $\sim 5\%$. This creates a seismic waveguide: energy injected into the LVZ undergoes total internal reflection at both boundaries where the impedance increases outward.

The trapping condition is:
\begin{equation}
    \sin\theta_c = \frac{V_{LVZ}}{V_{surround}} \approx \frac{4.35}{4.65} \approx 0.936 \quad \Longrightarrow \quad \theta_c \approx 69^\circ
\end{equation}

Seismic energy arriving at angles greater than $\theta_c$ is perfectly reflected back into the channel---structurally identical to optical fibre total internal reflection and the QCD flux-tube confinement derived in Volume~IV. The universal impedance operator governs wave trapping from $10^{-15}$ m (nucleon) to $10^{6}$ m (planetary mantle).

